\documentclass{article}
\usepackage[utf8]{inputenc}
\usepackage{amsmath, amssymb}
\usepackage{siunitx}
\usepackage{booktabs}
\title{Detaljerte beregninger for EMP-generator (10 J)}
\author{}
\date{}

\begin{document}
\maketitle

\textbf{Merk: Alle beregninger er teoretiske. Faktiske verdier avhenger av målte parametere som spoleinduktans, motstand og gnistgapkarakteristikk.}

\section{Hovedkondensator og spole}

\subsection{Spesifikasjoner}
\begin{itemize}
\item Kondensator: $C = 100\,\mu\text{F}$, merkespenning $450\,\text{V}$ (lades til $V_0 = 450\,\text{V}$)
\item Spole: Flat spiral, $N = 4$ vindinger, indre diameter $20\,\text{mm}$, ytre diameter $100\,\text{mm}$, tråd: $10\,\text{AWG}$ ($2.588\,\text{mm}$ diameter)
\end{itemize}

\subsection{Induktans for flat spiralspole (Wheelers formel)}
Indre radius $r_i = 10\,\text{mm}$, ytre radius $r_o = 50\,\text{mm}$.
Middelradius $r = (r_i + r_o)/2 = 30\,\text{mm} = 0.03\,\text{m}$.
Viklingsbredde $w = r_o - r_i = 40\,\text{mm} = 0.04\,\text{m}$.
I tommer: $r = 1.1811''$, $w = 1.5748''$.
\begin{equation}
L = \frac{r^2 N^2}{8r + 11w} = \frac{1.1811^2 \cdot 4^2}{8 \cdot 1.1811 + 11 \cdot 1.5748} = \frac{1.395 \cdot 16}{9.4488 + 17.3228} = \frac{22.32}{26.7716} \approx 0.8336\,\mu\text{H}.
\end{equation}
Vi bruker $L = 0.83\,\mu\text{H}$ som et estimat. Målt verdi kan avvike betydelig ($\pm 20\%$ eller mer).

\subsection{Energi lagret i kondensator}
\begin{equation}
E_0 = \frac{1}{2} C V_0^2 = 0.5 \cdot 100 \times 10^{-6} \cdot 450^2 = 10.125\,\text{J} \approx 10\,\text{J}.
\end{equation}

\section{Toppstrøm og pulsform}

\subsection{Karakteristisk impedans}
\begin{equation}
Z_0 = \sqrt{\frac{L}{C}} = \sqrt{\frac{0.83 \times 10^{-6}}{100 \times 10^{-6}}} = \sqrt{0.0083} \approx 0.0911\,\Omega.
\end{equation}

\subsection{Toppstrøm (ideell, tapsfri)}
\begin{equation}
I_{\text{peak, ideal}} = \frac{V_0}{Z_0} = \frac{450}{0.0911} \approx 4940\,\text{A}.
\end{equation}

\subsection{Dempet RLC-analyse}
Total motstand $R$ i kretsen er dominert av ESR i kondensator, ledningsmotstand og lysbuemotstand. Vi antar $R \approx 0.1\,\Omega$ (må verifiseres eksperimentelt). Dempingsfaktor:
\begin{equation}
\alpha = \frac{R}{2L} = \frac{0.1}{2 \cdot 0.83 \times 10^{-6}} \approx 60241\,\text{s}^{-1}.
\end{equation}
Resonansvinkelfrekvens:
\begin{equation}
\omega_0 = \frac{1}{\sqrt{LC}} = \frac{1}{\sqrt{0.83 \times 10^{-6} \cdot 100 \times 10^{-6}}} = \frac{1}{\sqrt{8.3 \times 10^{-11}}} \approx 109719\,\text{rad/s}.
\end{equation}
Dempet vinkelfrekvens:
\begin{equation}
\omega_d = \sqrt{\omega_0^2 - \alpha^2} = \sqrt{109719^2 - 60241^2} \approx 91700\,\text{rad/s}.
\end{equation}
Frekvens: $f_d = \omega_d/(2\pi) \approx 14.6\,\text{kHz}$.

Første topp inntreffer ved $t_p$ der $\tan(\omega_d t_p) = \omega_d/\alpha$:
\begin{equation}
\omega_d t_p = \arctan(91700/60241) \approx 0.99\,\text{rad} \Rightarrow t_p = 0.99/91700 \approx 10.8\,\mu\text{s}.
\end{equation}
Toppstrøm:
\begin{equation}
I_{\text{peak}} = \frac{V_0}{\omega_d L} e^{-\alpha t_p} \sin(\omega_d t_p) = \frac{450}{91700 \cdot 0.83\times10^{-6}} e^{-0.650} \cdot 0.836 \approx 5913 \cdot 0.522 \cdot 0.836 \approx 2580\,\text{A}.
\end{equation}
Med antatt $R=0.1\,\Omega$ er toppstrømmen $\sim 2600$ A. Denne verdien er følsom for $R$ og må måles.

\subsection{Pulsvarighet}
Periode: $T_d = 1/f_d \approx 68.5\,\mu\text{s}$. Tidskonstant for innhyllingskurven: $\tau = 1/\alpha \approx 16.6\,\mu\text{s}$. Etter $5\tau \approx 83\,\mu\text{s}$ er svingningene nesten borte.

\section{Magnetfelt}

\subsection{Felt i sentrum av spolen (sum over vindinger)}
Antar jevnt fordelte vindinger med radier: $r_1 = 10\,\text{mm}$, $r_2 = 23.33\,\text{mm}$, $r_3 = 36.67\,\text{mm}$, $r_4 = 50\,\text{mm}$.
\begin{equation}
B(0) = \frac{\mu_0 I}{2} \sum_{k=1}^{4} \frac{1}{r_k} = \frac{4\pi\times10^{-7} \cdot 2600}{2} \left( \frac{1}{0.01} + \frac{1}{0.02333} + \frac{1}{0.03667} + \frac{1}{0.05} \right).
\end{equation}
Summen av $1/r_k$: $100 + 42.86 + 27.27 + 20 = 190.13\,\text{m}^{-1}$.
\begin{equation}
B(0) = 1.633\times10^{-3} \cdot 190.13 \approx 0.31\,\text{T}.
\end{equation}
Dette er et estimat; faktisk felt kan avvike på grunn av spiralform og trådtykkelse.

\subsection{Feltets avtagning med avstand}
For $z = 0.2\,\text{m}$:
\begin{equation}
B(z) = \frac{\mu_0 I}{2} \sum_{k=1}^{4} \frac{r_k^2}{(r_k^2 + z^2)^{3/2}}.
\end{equation}
Beregner hvert ledd:
\begin{align}
r_1 &: \frac{0.01^2}{(0.01^2+0.2^2)^{1.5}} = \frac{1\times10^{-4}}{0.0401^{1.5}} \approx 0.01244, \\
r_2 &: \frac{0.02333^2}{0.04054^{1.5}} \approx 0.0667, \\
r_3 &: \frac{0.03667^2}{0.04134^{1.5}} \approx 0.160, \\
r_4 &: \frac{0.05^2}{0.0425^{1.5}} \approx 0.285.
\end{align}
Sum: $0.524$. $B(0.2) = 1.633\times10^{-3} \cdot 0.524 \approx 8.6\times10^{-4}\,\text{T} = 0.86\,\text{mT}$.

\section{Indusert spenning i pickup-spole}
Pickup-spole: $N_p = 5$, areal $A_p = 1\,\text{cm}^2 = 1\times10^{-4}\,\text{m}^2$, plassert på aksen i avstand $z = 0.2\,\text{m}$.
Maksimal $dB/dt$: $B(0.2) \cdot \omega_d = 8.6\times10^{-4} \cdot 91700 \approx 78.9\,\text{T/s}$.
\begin{equation}
\mathcal{E}_{\text{max}} = N
\mathcal{E}_{\text{max}} = N_p A_p \frac{dB}{dt} \approx 5 \cdot 1\times10^{-4} \cdot 78.9 \approx 0.0395\,\text{V} = 39.5\,\text{mV}.
\end{equation}
Dette er et estimat; faktisk verdi avhenger av nøyaktig geometri, orientering og strømpulsform.

\section{Komponentpåkjenninger}

\subsection{MOSFET i ladekrets}
Under avslag kan drain-spenningen nå $V_{\text{ds}} = V_{\text{bat}} + V_{\text{reflektert}} + V_{\text{leak}}$. Med $V_{\text{bat}} = 45\,\text{V}$, viklingsforhold $n = N_s/N_p = 200/10 = 20$, reflektert spenning $V_{\text{reflektert}} = V_{\text{ut}}/n = 450/20 = 22.5\,\text{V}$. Lekkinduktansspike kan være betydelig; snubberen demper den, men uten måling kan vi ikke angi eksakt verdi. IRFP450 tåler $500\,\text{V}$, så marginen er sannsynligvis tilstrekkelig, men må verifiseres.

\subsection{Dioder}
UF4007 har $1000\,\text{V}$ sperrespenning. For flyback-utgangen er reversspenningen $V_{\text{rev}} = V_{\text{ut}} + V_{\text{bat}}/n = 450 + 45/20 = 452.25\,\text{V}$ pluss transienter. Én UF4007 er marginal; vi anbefaler to i serie med spenningsdelingsmotstander ($10\,\text{M}\Omega$ over hver) for å oppnå $>1500\,\text{V}$ rating.

\subsection{Gnistgap}
Wolframelektroder har et smeltepunkt på $3422\,^{\circ}\text{C}$ og tåler høye pulsstrømmer. Avstanden på $0.5\,\text{mm}$ gir en naturlig gjennomslagsspenning på ca. $1500\,\text{V}$ i tørr luft, men med trigging kan den utløses ved lavere spenning. Faktisk triggespenning må karakteriseres eksperimentelt.

\subsection{Snubber-komponenter}
R2 og C2 dimensjoneres for å absorbere lekkinduktansenergien. Med antatt lekkinduktans $L_{\text{leak}} \approx 5\%$ av $L_p$ ($130\,\mu\text{H} \rightarrow 6.5\,\mu\text{H}$) og toppstrøm $\sim 2\,\text{A}$, er energien $\frac{1}{2} L_{\text{leak}} I^2 = 0.5 \cdot 6.5\times10^{-6} \cdot 4 = 13\,\mu\text{J}$. En $10\,\text{nF}$ kondensator og $10\,\text{k}\Omega$ motstand gir en tidskonstant på $100\,\mu\text{s}$, som er tilstrekkelig.

\section{Termiske betraktninger (hypotetiske)}
Energifordelingen er ukjent. Eksempel: Hvis $20\%$ av $10\,\text{J}$ ($2\,\text{J}$) går til oppvarming av spolen (masse $\sim 89\,\text{g}$, varmekapasitet $385\,\text{J/(kg·K)}$):
\begin{equation}
\Delta T \approx \frac{2}{0.089 \cdot 385} \approx 0.058\,\text{K}.
\end{equation}
For kondensator (masse $\sim 30\,\text{g}$, varmekapasitet $\sim 30\,\text{J/K}$) med $30\%$ tap ($3\,\text{J}$):
\begin{equation}
\Delta T \approx \frac{3}{30} = 0.1\,\text{K}.
\end{equation}
Lokale varmepunkter kan likevel oppstå. Gjentatte pulser krever kjøling.

\section{Repetisjonsrate}
Ladestrømmen fra flyback-omformeren er i gjennomsnitt ca. $20$--$50\,\text{mA}$. For å lade $100\,\mu\text{F}$ til $450\,\text{V}$ trengs $Q = C \cdot V = 0.045\,\text{C}$. Med $30\,\text{mA}$ tar det $t = Q/I = 0.045/0.03 = 1.5\,\text{sekunder}$. I praksis vil strømmen avta etter hvert som spenningen stiger, så ladetid $2$--$5\,\text{sekunder}$ er realistisk. Repetisjonsraten blir da $0.2$--$0.5\,\text{Hz}$. Med alkaliske 9V-batterier er dette maksimalt; LiPo-batterier kan gi høyere strøm og raskere lading.

\section{Konklusjon}
Alle beregninger er teoretiske og avhenger av antatte parametere. Fysisk prototype må bygges og måles for å bekrefte ytelse og sikkerhet.
\end{document}